% BT1899 -- combined Holonet insertion: residual-language clarification + guard-envelope theorem
% Suggested insertion points:
%   (1) residual clarification: Appendix B, after architecture-completeness opening paragraphs.
%   (2) guard-envelope theorem: near the Witting transaction / Fano time-bin envelope discussion.

\subsection*{BT1899 Residual and Guard-Envelope Clarification}

\paragraph{Finite closure versus physical-continuum closure.}
The Holonet architecture is complete as a finite machine specification: the $W(3,3)$ carriers,
routing atlas, Steinberg memory, Witting transaction ABI, contextual-fuel accounting, and finite
lattice/register data are exact or machine-witnessed.  This finite closure must be kept distinct from
physical-continuum closure.  The safe public statement is
\[
\text{finite machine-complete architecture, with remaining physical/continuum identifications classified.}
\]
In this language, R1 fixes finite $E_8/2E_8$ and positive-$E_8$ gauge-lattice data but does not by
itself complete the Einstein--Hilbert continuum derivation.  R2 supplies moonshine traces and the
finite matter-register bridge while leaving physical fermion labelling and full CKM/PMNS derivations
as separate tasks unless proved elsewhere.  R3 gives a theorem-governed convergence route, not yet a
complete constructive physical continuum coefficient.

\paragraph{Fano time-bin guard envelope.}
The Witting transaction table has $1600$ active frames.  Embedding these frames in an eleven-bit
single-photon time-bin register gives
\[
2^{11}=2048,
\qquad
2048-1600=448=7\cdot64.
\]
The slack is naturally organized as seven Fano guard pages.  Each page has
\[
64=24+24+16
\]
guard bins, interpreted as
\[
24\text{ dark-reference bins}
+24\text{ loss-probe bins}
+16\text{ parity-overflow bins}.
\]
Thus the guard envelope is a repo-derived BT1649--BT1650 theorem and should be cited as such unless
this insert is included in the paper body.
